\chapter{Conclusiones y Trabajo Futuro\label{conclusiones}}

En este último capítulo se recopilan las conclusiones finales extraídas a lo largo del diseño, desarrollo e implementación de la plataforma inmersiva de auditoría arquitectónica. Se evalúa el grado de cumplimiento de los objetivos establecidos al inicio del proyecto, se analiza el impacto potencial del sistema en el sector de la edificación y, por último, se proponen las líneas de trabajo futuro para la evolución del software más allá del marco de este Trabajo de Fin de Grado.

\section{Consecución de Objetivos}
Tras finalizar el desarrollo de la herramienta y la realización de las pruebas en el motor gráfico, se puede afirmar que los objetivos propuestos originalmente se han alcanzado en su mayoría. El propósito principal de reducir la brecha entre el diseño teórico de un edificio (BIM) y su ejecución real en obra se ha resuelto con éxito mediante la fusión de tecnologías XR y reconstrucción volumétrica tridimensional (\textit{Gaussian Splatting}).

Específicamente, se ha consolidado un entorno inmersivo estable y fluido capaz de procesar nubes de puntos complejas previamente optimizadas mediante técnicas de submuestreo espacial y \gls{sor}. Las herramientas operacionales como el escáner de desviaciones mediante mapas de calor, la cinta métrica tridimensional, el detector automatizado de colisiones estructurales con aislamiento de componentes y el módulo de captura fotográfica se encuentran plenamente integradas y operativas. La aplicación no solo actúa como un visualizador, sino que cumple con el requisito de unificar todos los datos en un motor procedural que compila de forma automática informes documentales en formato PDF listos para su uso profesional.

\section{Repercusión en el Entorno Destino}
La introducción de soluciones de Realidad Virtual orientadas al flujo de trabajo de construcción influencia positivamente a la comunidad y la industria de la construcción:
\begin{itemize}
    \item \textbf{Reducción de costes por errores de ejecución:} En la gestión de obras tradicional, un error geométrico o una interferencia entre una tubería MEP y un elemento estructural que no se detecta a tiempo se traduce en demoras masivas y sobrecostes económicos de demolición y reconstrucción. Esta herramienta permite anticipar dichos conflictos en fases tempranas mediante una inspección visual intuitiva a escala real 1:1.
    \item \textbf{Democratización del acceso a datos técnicos:} Tradicionalmente, la interpretación de nubes de puntos densas o modelos de coordinación requería software de escritorio complejo y personal altamente cualificado. Al trasladar esta información a una experiencia inmersiva natural, cualquier agente involucrado (promotores, arquitectos, contratistas o clientes finales) puede comprender las medidas y el estado real de la obra sin necesidad de formación técnica avanzada.
    \item \textbf{Optimización del flujo de información:} La capacidad de generar informes técnicos unificados directamente desde el visor VR acorta los tiempos de comunicación entre el personal de campo y la oficina técnica, acelerando la toma de decisiones y el control de calidad.
\end{itemize}

\section{Limitaciones y Líneas de Trabajo Futuro}
A pesar de haber alcanzado un producto mínimo viable completamente funcional que cumple con los objetivos principales del proyecto, el proceso de desarrollo e investigación ha demostrado ciertas limitaciones técnicas y ventanas de oportunidad. Debido a las restricciones temporales inherentes a un Trabajo de Fin de Grado y a la complejidad de las tecnologías XR, las siguientes funcionalidades no pudieron ser integradas en la versión actual y se proponen como futuras líneas de investigación y desarrollo:

\begin{itemize}
    \item \textbf{Integración de Realidad Aumentada (AR):} Aunque la aplicación resuelve de manera óptima la inspección inmersiva en entornos virtuales (VR), quedó fuera del alcance la implementación de un módulo de Realidad Aumentada. Esta funcionalidad habría permitido superponer el modelo digital BIM directamente sobre el terreno físico real de la obra utilizando técnicas de geolocalización y visión artificial, lo cual potenciaría las capacidades de auditoría sobre el propio emplazamiento.
    \item \textbf{Automatización de la importación:} Actualmente, la incorporación de un nuevo edificio requiere una serie de procesos de configuración manual dentro del editor de Unity (organización jerárquica de capas, preparación de los colisionadores y ajuste de variables del escáner). Una línea futura clave consiste en el desarrollo de un algoritmo en tiempo de ejecución que automatice por completo la creación de la escena con solo depositar el archivo del modelo, eliminando la dependencia del entorno de desarrollo.
    \item \textbf{Entorno colaborativo multiusuario:} La arquitectura actual de la aplicación está diseñada para sesiones individuales de un único usuario. La integración de un sistema de red y sincronización (como \textit{Unity Netcode} o \textit{Photon Fusion}) permitiría habilitar sesiones multiusuario concurrentes, facilitando que múltiples inspectores, arquitectos y directores de obra interactúen simultáneamente en la misma maqueta virtual para discutir las incidencias de forma remota.
    \item \textbf{Fidelidad y estética de la Interfaz de Usuario (UI):} Los paneles y menús desarrollados cumplen estrictamente con los requisitos funcionales de diseño responsivo, técnicas anti-clipping y prevención de fatiga visual. No obstante, a nivel estético presentan un acabado básico e industrial. Un desarrollo posterior enfocado en el pulido estético permitiría generar interfaces más detalladas, integrando elementos interactivos avanzados y menús que eleven la calidad visual del producto.
    \item \textbf{Sincronización en la nube y conectividad:} El motor de informes actual exporta los documentos PDF y las capturas fotográficas guardándolos de forma local en el almacenamiento persistente del visor Meta Quest. Una mejora a considerar sería la implementación de un servicio de sincronización automática con repositorios en la nube (como \textit{Firebase} o servidores BIM corporativos), permitiendo que los informes de auditoría estén disponibles de manera instantánea en las oficinas tras concluir la inspección.
    \item \textbf{Recorrido virtual guiado (tour automatizado):} Inicialmente se contempló la inclusión de un sistema de vuelo cinemático "sobre raíles" para presentar el modelo BIM de forma automática. Sin embargo, esta funcionalidad fue descartada durante la fase de desarrollo debido a que forzar el desplazamiento y la rotación de la cámara generaba cinetosis con excesiva facilidad. Sumado a los conflictos técnicos que presentaba forzar ese movimiento de la cámara y que se trataba de una característica secundaria respecto a las herramientas centrales de auditoría, se decidió prescindir de ella para priorizar la estabilidad del entorno.
\end{itemize}

%\section{Dificultad Global del Proyecto e Impresiones Personales}
%El desarrollo de este Trabajo de Fin de Grado ha supuesto dentro del grado de Ingeniería Multimedia un desafío de gran envergadura. La principal dificultad técnica no solo fue la programación aislada de las mecánicas de interacción, sino en la convergencia de estándares industriales dispares (formatos IFC provenientes de la arquitectura frente a mallas poligonales optimizadas para videojuegos) y en las severas restricciones de rendimiento gráfico que imponen los visores de Realidad Virtual autónomos (\textit{Standalone}). 

%Equilibrar la balanza entre la alta fidelidad visual que exige una auditoría de obra y el mantenimiento de una tasa de fotogramas por segundo (FPS) elevada para garantizar el confort del usuario supuso un exhaustivo proceso de optimización arquitectónica. Fue necesario abordar problemas complejos de computación gráfica, tales como la evasión de cuellos de botella en la memoria RAM mediante corrutinas asíncronas para la asignación procedimental de colisionadores, la simplificación del cálculo de distancias euclidianas eludiendo raíces cuadradas pesadas en la CPU y la reconfiguración del búfer gráfico para el renderizado estable de mapas de calor.

%A nivel personal, este proyecto ha constituido una de las experiencias de aprendizaje más enriquecedoras de toda la etapa universitaria. Ha demandado una actitud proactiva para investigar tecnologías emergentes en constante evolución (como el \textit{Gaussian Splatting}) y ha forzado la aplicación transversal de disciplinas del grado, desde la algoritmia pura y las estructuras de datos hasta el diseño de interfaces inmersivos. La satisfacción de transformar líneas de código abstractas en una herramienta tangible y de alta utilidad industrial reafirma la vocación ingenieril del autor.